% =============================================================================
%  DRAFT manuscript -- the circular-evaluation paper (reframed from TopoBrain).
%  New file; does NOT touch paper_iet/paper.tex (the old, refuted paper).
%  Venue-neutral article class for now (reformat to Sci Reports / MELBA / IEEE later).
%  Numbers marked [PENDING] await the confound-free eval and the cross-method baselines.
%  Everything else is the real, verified n=10 result.
% =============================================================================
\documentclass[11pt,a4paper]{article}
\usepackage[utf8]{inputenc}
\usepackage[T1]{fontenc}
\usepackage[margin=25mm]{geometry}
\usepackage{amsmath,amssymb}
\usepackage{booktabs}
\usepackage{graphicx}
\graphicspath{{figures/}}
\usepackage[hidelinks]{hyperref}
\usepackage{lineno}\linenumbers

\title{\bfseries Self-Confirming Fidelity: A Generator-Coupled Segmenter Masks
Over-Smoothing that an Independent Segmenter Reveals in MRI Synthesis Evaluation}
\author{P.\ Bashyal, P.\ Adhikari, S.\ Bhandari, N.\ B.\ Adhikari\\
\small Department of Electronics and Computer Engineering, Pulchowk Campus, IOE, Tribhuvan University}
\date{}

\begin{document}
\maketitle

\begin{abstract}
\noindent Structural-fidelity claims for synthesised MRI are increasingly supported by
segmentation-based evaluation: a segmentation network judges whether synthetic anatomy is
preserved. We show that this practice is \emph{circular} when the segmenter is coupled to the
generator, and that the circularity is large enough to invert the conclusion. Using paired
3\,T/7\,T brain MRI (10 subjects) and a deterministic 3T$\rightarrow$7T translation network with an
auxiliary segmentation head, we measure the topology (connected-component count, $\beta_0$) of
synthetic grey matter under two judges applied to the \emph{identical} images: the model's own
co-trained head, and an independent segmenter that never saw the generator. The co-trained head
rates synthetic grey matter as structurally comparable to---or richer than---real 7\,T (median
synthetic/real $\beta_0$ ratio $1.38$), whereas the independent segmenter, on the same images,
shows it is roughly ten times topologically simpler---severe over-smoothing (median ratio $0.10$;
$0.09$--$0.12$ in every subject). In all 10 subjects the coupled judge rates grey-matter topology
$5$--$31\times$ higher than the independent judge, including both held-out subjects, so the reversal
is a property of the metric, not of overfitting. The bias is a dose-response in coupling: with all
three segmenters applied to the \emph{same} images, certification of fidelity weakens monotonically
as the segmenter is decoupled from the generator ($1.38 \rightarrow 0.95 \rightarrow 0.10$). We further show that standard
whole-volume PSNR/SSIM overstate fidelity by ${\sim}8$\,dB / $0.45$ SSIM relative to brain-masked
metrics. We conclude that only an independent, real-data-trained segmenter is an admissible
structural judge for synthesis, and release the evaluation protocol.
\end{abstract}

\section{Introduction}
Cross-field-strength MRI synthesis (e.g.\ 3T$\rightarrow$7T) is evaluated less and less by
intensity metrics alone, because PSNR/SSIM are known to miss clinically relevant morphological
error. A common response is to evaluate \emph{structural} fidelity with a segmentation network:
segment the synthetic image, and check that its anatomy matches the target. When the segmenter is
\emph{coupled} to the generator---jointly trained, or sharing features or data---this evaluation
can confirm itself: the metric was optimised to like exactly the outputs it is now asked to judge.

We show this failure mode is real, and large. On paired 3T/7T brain MRI, a segmentation head
co-trained with a 3T$\rightarrow$7T generator certifies that synthetic grey-matter topology is
preserved; an independent segmenter applied to the same synthetic images shows the opposite---the
grey matter is roughly ten times topologically simpler than real 7\,T, the signature of
over-smoothing. The verdict \emph{reverses} depending only on which segmenter judges it, in 9 of
10 subjects. We further show the effect scales with coupling.

The contribution is a validation-methodology finding, not a synthesis method; the translation
network is a demonstration vehicle. Our contributions are:
\begin{enumerate}
\item The first empirical demonstration that a generator-coupled segmenter yields a
self-confirming structural-fidelity verdict for medical image synthesis, and that an independent
segmenter reverses it, with a consistent, quantified effect (Section~\ref{sec:reversal}).
\item Evidence that the bias is a \emph{dose-response} in coupling: joint $0.98\rightarrow$
detached $0.41\rightarrow$ independent $0.10$ (Section~\ref{sec:dose}).
\item A protocol prescription---only an independent, real-data-trained segmenter is an admissible
structural judge---together with a corroborating account of whole-volume metric inflation
(Section~\ref{sec:inflation}).
\end{enumerate}

\section{Related work}
This work extends the evaluation-pitfalls literature---exemplified by \emph{Metrics Reloaded}
\cite{maierhein2024metricsreloaded} and the ``common limitations of image-processing metrics''
programme---from segmentation/detection to \emph{generative synthesis}. Deo et al.\
\cite{deo2025metrics} show that no-reference image-quality metrics miss localised morphological
changes and advocate downstream-task evaluation; Dohmen et al.\ \cite{dohmen2025similarity}
analyse similarity/quality metrics for MR image-to-image translation. Neither contrasts a
co-trained against an independent judge, nor reports a verdict reversal. In 3T$\rightarrow$7T
synthesis specifically, Acs and Zhuang \cite{acs2025semisupervised} evaluate with an
\emph{independent} segmentation pipeline---thereby avoiding the bias we identify without naming or
measuring it. Coupled segmenters are, meanwhile, routinely built \emph{into} generator training
(shape- and cycle-consistency losses \cite{zhang2018translating}); we show that grading such models
with their own segmenter is circular. Topology in generative MRI has been used as a
generation-time prior (e.g.\ persistence-guided diffusion \cite{bhattacharya2025brainmrdiff}),
not as an evaluation diagnostic of output complexity.

\section{Materials and methods}
\paragraph{Data.} Ten healthy adults with paired 3\,T and 7\,T T1-weighted volumes
(0.65\,mm isotropic), resampled to a common grid. A complete FreeSurfer \texttt{aparc+aseg}
parcellation of each real 7\,T volume is mapped to a 4-class tissue label
(background, CSF, grey matter (GM), white matter (WM)) and used only as \emph{training}
supervision for the auxiliary head---not as a topology reference (see Section~\ref{sec:reversal}).

\paragraph{Model (a vehicle, not a contribution).} A deterministic cascaded translation network
maps a 3\,T volume to a synthetic 7\,T image via a 3D U-Net generator, and an auxiliary
segmentation head reads the generated image to predict tissue class. The head is trained jointly
with the generator; because its input is the generator's output, it is \emph{coupled} to the
generator. Training is deterministic and seeded; the same model is used throughout.

\paragraph{Two-judge evaluation.} For each subject we obtain the synthetic 7\,T image and score its
grey-matter topology under two segmenters applied to the identical image within an identical brain
mask: (i) the model's own \emph{co-trained head}; (ii) an \emph{independent} judge---an
unsupervised three-class Gaussian-mixture tissue segmentation of the image intensities, which never
saw our training data, reads only the image (no conditioning), and---crucially---enforces no
topology prior. We report per-tissue $\beta_0$ (the number of 26-connected components), and the
synthetic-vs-real ratio $r = \beta_0(\text{synthetic})/\beta_0(\text{real})$: $r\!\approx\!1$ reads
as ``structurally comparable / preserved''; $r\!\ll\!1$ reads as ``over-smoothed'' (the synthetic
tissue is far simpler than real). Comparisons are made only \emph{within} a judge, never across
judges, since absolute $\beta_0$ is segmenter- and threshold-dependent.

\paragraph{Why FreeSurfer is not the reference.} A FreeSurfer $\beta_0$ is an algorithmic guarantee
of its topology-corrected surface reconstruction, not a measurement of the image; comparing any
voxel-wise segmenter to it is confounded. Hence the within-judge synthetic-vs-real design.

\paragraph{Dose-response.} To test whether the bias tracks coupling, we additionally train a
\emph{detached} variant---identical in every respect (same split, iterations, losses) except that
the segmentation head's gradient is stopped before the generator, so the head is trained on the
generator's outputs but does not shape them. This gives a coupling axis:
joint (gradient-coupled) $\rightarrow$ detached (output-coupled) $\rightarrow$ independent
(uncoupled).

\section{Results}
\subsection{The verdict reverses with the choice of segmenter}\label{sec:reversal}
Table~\ref{tab:reversal} reports the grey-matter synthetic/real $\beta_0$ ratio under each judge,
all applied to identical images across the 10 subjects. The co-trained head certifies synthetic
grey matter as structurally comparable to, or richer than, real 7\,T (median $r=1.38$). The
independent segmenter, on the same images, shows it is roughly ten times topologically simpler
(median $r=0.10$; per-subject $0.09$--$0.12$, i.e.\ tight). In every one of the 10 subjects the
co-trained judge rates grey-matter topology $5$--$31\times$ higher than the independent judge:
the two disagree on whether the synthesis preserved anatomy in all cases
(Fig.~\ref{fig:montage} visualises this for one subject; Fig.~\ref{fig:reversal}, all ten). The two \emph{held-out}
subjects behave identically to the eight in-sample ones, so the reversal is a property of the
metric, not of the model having memorised the training subjects. The effect is specific to grey
matter, the most structurally complex tissue; for CSF and WM both judges agree the synthesis is
over-smoothed (no reversal), which we report as-is.

\begin{figure}[t]\centering
\includegraphics[width=\linewidth]{fig_montage_sub-06.pdf}
\caption{The same synthetic image, judged by two segmenters (subject sub-06, one axial slice).
Grey-matter connected components are individually coloured, so the number of colours \emph{is}
$\beta_0$. The co-trained head yields a few large components ($\beta_0=134$, ``looks
preserved''); an independent segmenter reveals the synthesis is shattered into many specks
($\beta_0=698$, over-smoothed). The verdict depends entirely on who is asked.}
\label{fig:montage}
\end{figure}

\begin{figure}[t]\centering
\includegraphics[width=0.92\linewidth]{fig_reversal.pdf}
\caption{Per-subject grey-matter synthetic/real $\beta_0$ ratio under the co-trained head (blue,
$\gtrsim 1$: ``preserved / richer than real'') versus an independent segmenter (orange,
$\approx 0.1$: $\sim$10$\times$ over-smoothed), on identical images. The coupled judge exceeds the
independent by $5$--$31\times$ in all 10 subjects.}
\label{fig:reversal}
\end{figure}

\begin{table}[t]\centering
\caption{Grey-matter synthetic/real $\beta_0$ ratio by judge, identical images, 10 subjects.
Co-trained head $\ge 1$ (``preserved / richer than real''); independent segmenter $\approx 0.1$
(``$\sim$10$\times$ over-smoothed''). The coupled judge exceeds the independent by $5$--$31\times$
in every subject.}\label{tab:reversal}
\begin{tabular}{lcc}
\toprule
Judge & median ratio & reads as \\
\midrule
Co-trained head (coupled)     & 1.38 & preserved / richer than real \\
Independent GMM (uncoupled)   & 0.10 & $\sim$10$\times$ over-smoothed \\
\bottomrule
\end{tabular}
\end{table}

\subsection{The bias scales with coupling}\label{sec:dose}
To test whether the certification tracks coupling---and to rule out that generator differences drive
it---we apply all three segmenters to the \emph{same} images (one generator's synthetic 7\,T and the
real 7\,T). Any difference is then attributable to the segmenter's coupling alone.
Table~\ref{tab:dose} shows the grey-matter synthetic/real $\beta_0$ ratio (median over all 10
subjects) falling monotonically as the segmenter is decoupled from the generator: joint
(gradient-coupled) $1.38 \rightarrow$ detached (output-coupled) $0.95 \rightarrow$ independent
(uncoupled) $0.10$. On identical images, the certification is therefore a dose-response in coupling
(Fig.~\ref{fig:dose}).

Two levels are visible, of very different size. Any segmenter trained on the generator's outputs
---joint \emph{or} detached---reports synthetic grey matter as comparable to, or richer than, real
($0.95$--$1.38$); only the fully independent segmenter reveals the over-smoothing ($0.10$). Being
trained on the generator's outputs at all is thus the \emph{dominant} leakage---a large step
present in every subject---whereas gradient coupling adds only a smaller median increment
($1.38$ vs.\ $0.95$) that is not consistent across every subject. Notably, the coupled head is not
merely optimistic: it reports the over-smoothed synthesis as \emph{more} structurally complex than
the real image it fails to reproduce.

\begin{figure}[t]\centering
\includegraphics[width=0.82\linewidth]{fig_dose_response.pdf}
\caption{Dose-response: grey-matter synthetic/real $\beta_0$ ratio as the segmenter is decoupled
from the generator, all judges applied to the \emph{same} images. Faint lines: per subject; bold:
median ($1.38 \rightarrow 0.96 \rightarrow 0.10$). The large, universal step is to full
independence.}
\label{fig:dose}
\end{figure}

\begin{table}[t]\centering
\caption{Grey-matter synthetic/real $\beta_0$ ratio (median, 10 subjects) with all three segmenters
applied to the \emph{same} images, vs.\ segmenter--generator coupling. Monotonic: the bias is the
segmenter's coupling, not the generator.}\label{tab:dose}
\begin{tabular}{lcc}
\toprule
Coupling (all judges, identical images) & median ratio & reads as \\
\midrule
Joint head (gradient-coupled)   & 1.38 & ``richer than real'' \\
Detached head (output-coupled)  & 0.95 & preserved \\
Independent GMM (uncoupled)     & 0.10 & $\sim$10$\times$ over-smoothed \\
\bottomrule
\end{tabular}
\end{table}

\subsection{Generality across methods}\label{sec:generality}
\textbf{[PENDING]} To show the pitfall is not specific to one architecture, we reproduce two
published synthesis models (FR-U-Net; pix2pix) on the same split and repeat the two-judge
evaluation, testing whether the method ranking depends on which segmenter judges it.

\subsection{Whole-volume metrics overstate fidelity}\label{sec:inflation}
As a second, well-known-but-routinely-ignored pitfall in the same lineage, brain-masked SSIM/PSNR
are substantially lower than the whole-volume values that dominate the literature
(brain-masked ${\sim}0.45$ SSIM / ${\sim}14.3$\,dB vs.\ whole-volume ${\sim}0.90$ / ${\sim}21.5$\,dB;
a ${\sim}0.45$ SSIM / ${\sim}8$\,dB inflation driven by the large, easy background), reproducing
inflated headline figures reported in prior work.

\section{Discussion}
The finding generalises to a simple principle---a \emph{leakage law} for synthesis evaluation:
any learned structural-fidelity metric that shares data, features, or joint training with the
generator measures the generator's own bias, not anatomy. It is the generative-evaluation analogue
of train/test leakage. The practical prescription follows directly: structural fidelity of
synthesised images must be judged by an \emph{independent}, real-data-trained segmenter, and
published results that used a coupled or co-trained segmenter for structural/segmentation-consistency
evaluation should be re-examined. That the reversal is specific to grey matter is mechanistically
consistent: a coupled segmenter can only mask a defect where the structure is complex enough to
hide it; for simpler tissues even the coupled judge reports the over-smoothing.

\section{Limitations}
The cohort is small ($n=10$), healthy, and single-site; we therefore claim the \emph{existence and
mechanism} of the failure mode, not universal magnitudes. Eight of ten subjects are in-sample for
the trained model; this is acceptable because the claim is about the metric, and the two held-out
subjects behave identically. The reversal is grey-matter-specific. The independent judge is an
unsupervised intensity segmenter (no topology prior); a standard tool (FastSurfer/SynthSeg) as a
second independent judge is left to strengthen the claim. $\beta_0$ is segmenter- and
threshold-dependent, so we compare only within a judge. We make no clinical or diagnostic claim.

\section{Conclusion}
Whether synthetic MRI ``preserves anatomy'' can depend entirely on who is asked. A segmenter
coupled to the generator certifies fidelity that an independent segmenter refutes, and the bias
scales with coupling. Structural-fidelity evaluation of synthesis is only trustworthy when the
judge is independent of the generator.

\bibliographystyle{plain}
% \bibliography{circularity}   % TODO: assemble .bib -- keys: maierhein2024metricsreloaded,
% deo2025metrics, dohmen2025similarity, acs2025semisupervised, zhang2018translating,
% bhattacharya2025brainmrdiff. NOTE: the paper previously miscited "Onofrey" -> correct is Dohmen.
\end{document}
